\chapter{Introducción}\label{chap:Intro}

Este documento corresponde al template de \LaTeX{} con la estructura básica mínima requerida para el desarrollo de una tesis de postgrado del departamento. Contiene comentarios, sugerencias y recomendaciones basadas en las reglas y experiencia del comité de postgrado del \gls{die}.

Para hacer uso de este template, se recomienda leerlo íntegramente y en detalle antes de comenzar a completar cada sección. Cada capítulo contiene comentarios y recomendaciones dedicados a esa sección en particular. Sin embargo, algunos de ellos aplican de manera transversal a todo el documento. Una vez leído, se puede proceder a reemplazar el contenido con el material propio de la tesis.

Entre los aspectos transversales más relevantes se encuentran la extensión del documento, que se sugiere sea de entre 80 y 120 páginas, y el manejo de la bibliografía. Sobre esta última, se recomienda utilizar el formato \gls{ieee}, pero en cualquier caso, se debe mantener un formato consistente. El formato \textit{ieeetr} de BibTeX incluido en este template facilita esta tarea; sin embargo, es importante asegurarse de lo siguiente:

\begin{itemize}
    \item Incluir los nombres de los autores con el formato correcto (Inicial. Apellido).
    \item Evitar abreviaturas.
    \item Al utilizar siglas, asegurarse de que estén definidas.
    \item Al citar sitios web, indicar la fecha en que fue consultada la información.
    \item Incluir los títulos completos; en inglés con mayúscula inicial en cada palabra, en español solo en la primera.
    \item En una tesis en español, los nombres de los autores deben ir separados por ``y''.
    \item Utilizar correctamente el tipo de citación (libro, revista, conferencia, \textit{tech report}, etc.).
    \item En referencias con muchos autores, nombrar algunos y utilizar \textit{et al.} para el resto.
    \item Las referencias de arXiv corresponden a trabajos aún no publicados formalmente; si el trabajo fue publicado posteriormente, citar la conferencia o revista correspondiente.
    \item Al referenciar instituciones u organizaciones, no se debe aplicar el formato de autor personal. Por ejemplo, el Coordinador Eléctrico Nacional no debe abreviarse como C. E. Nacional, sino citarse con su nombre completo. Para ello se puede usar  \texttt{\{\{Coordinador Eléctrico Nacional}\}\} en BibTeX para evitar que se interprete como nombre y apellido.
\end{itemize}

Todas las referencias incluidas en el documento deben ser citadas en el texto; no se deben incluir referencias que no sean mencionadas.

Sobre este capítulo en particular, en él se debe presentar el contexto general de la tesis, incluyendo la motivación, la hipótesis de investigación y los objetivos del trabajo. No se debe incluir una revisión del estado del arte. En esta plantilla se incluye la estructura básica, pero se pueden agregar secciones adicionales, como contribuciones o alcances, si el autor lo considera pertinente.

\section{Motivación}

En esta sección se debe describir por qué se está investigando el tema de la tesis. Una manera intuitiva de lograrlo es introducir el tema desde lo más general hasta lo más específico, señalando su importancia en el área. Luego, se debe identificar el problema a resolver en ese ámbito específico, mencionar las soluciones típicas existentes y contrastar con la solución novedosa que se propone en la tesis. Para esta sección normalmente se incluirán afirmaciones que deben estar respaldadas por una o varias referencias. 
\section{Hipótesis}
A continuación, se debe detallar la hipótesis de investigación. Se sugiere usar una sola hipótesis en el caso de magíster. Es importante que esta hipótesis sea validable, ya que debe demostrarse su cumplimiento a partir de los resultados. Para ello, es recomendable incluir alguna métrica que ayude a verificar dicho cumplimiento, ya sea de forma directa o indirecta. Por ejemplo: ``(...) consiguiendo una tasa de error de bit no superior a $10^{-3}$''.

\section{Objetivos}
En esta sección se deben detallar los objetivos a lograr en la tesis, los cuales deben ser claros y medibles, ya que su cumplimiento debe ser demostrado y evidenciado a lo largo de las siguientes secciones. Para ello, se recomienda utilizar frases que empiecen con verbos en infinitivo; una referencia útil para la selección de estos verbos es la taxonomía de Bloom \cite{bloom_taxonomy}, que clasifica los verbos según el nivel cognitivo de la acción. Se debe evitar el uso de verbos ambiguos como ``estudiar'', los cuales dan a entender que no se realiza nada nuevo y son difícilmente medibles.
En cuanto a la estructura, se recomienda incluir un objetivo general y un máximo de cinco objetivos específicos. El objetivo general debe estar directamente relacionado con el título de la tesis, mientras que los objetivos específicos deben ser consistentes con la hipótesis de investigación planteada anteriormente.
\subsection{Objetivo General}
Desarrollar y evaluar...
\subsection{Objetivos Específicos}
\begin{itemize}
    \item Caracterizar...
    \item Diseñar...
    \item Aplicar...
    \item Validar...
\end{itemize}
\section{Estructura de la tesis}
En esta sección se describen el resto de capítulos de la tesis. Por ejemplo: ``El resto de este documento está organizado de la siguiente manera: El capítulo \ref{chap:MarcoT} introduce los conceptos fundamentales de (...). El capítulo \ref{chap:metodologia} introduce la metodología utilizada para (...). El capítulo \ref{chap:resultados} analiza los resultados de (...). Finalmente, el capítulo \ref{chap:conclusion} discute el cumplimiento de los objetivos, señalando posibles espacios de mejora...''